\pdfoutput=1
\documentclass[12pt]{article}

\usepackage{amsmath, amssymb, amsthm}
\usepackage{graphicx}
\usepackage{booktabs}
\usepackage{array}
\usepackage{longtable}
\usepackage{calc}
\usepackage{etoolbox}
\makeatletter
\patchcmd\longtable{\par}{\if@noskipsec\mbox{}\fi\par}{}{}
\makeatother
\IfFileExists{footnotehyper.sty}{\usepackage{footnotehyper}}{\usepackage{footnote}}
\makesavenoteenv{longtable}
\usepackage[colorlinks=true, linkcolor=blue, citecolor=blue, urlcolor=blue]{hyperref}
\usepackage{geometry}
\geometry{margin=1in}

\providecommand{\tightlist}{%
  \setlength{\itemsep}{0pt}\setlength{\parskip}{0pt}}

\title{The Atom as a Measurable Tholon:\\
\large From Einstein's Incomplete Electron to a Field-First Ontology}
\author{J. W. Milton\\[4pt]{\small Clarity Coalition}}
\date{10 April 2026 \\ \small v1.0}

\begin{document}
\maketitle

\noindent\textbf{Provisional arXiv subjects:} hep-ph (secondary: physics.hist-ph; math-ph)

\begin{abstract}
The tholonic model proposes that all stable, self-sustaining structures
in nature are instances of a triadic pattern called the N-D-C triad
(Negotiation, Definition, Contribution). This paper argues that the atom
is not merely analogous to a tholon but is measurably one: several of
its quantitatively verified properties --- quark occupancy structure,
the virial theorem's kinetic/potential energy balance, color charge
structure, tetrahedral orbital geometry, and the Koide formula's
unexplained \(2/3\) lepton mass ratio --- are structural predictions of
the tholonic framework rather than independent empirical accidents. The
Tsirelson bound (\(2\sqrt{2}\)) in Bell inequality experiments is
derived from the real-virtual tholon pair structure at two independent
levels: (a) the complete tholon's composite convergence
(\(\sqrt{2} + \sqrt{2} = 2\sqrt{2}\), each component proved in Paper 1);
and (b) pure angle geometry, where the CHSH maximum equals
\(4\cos(\pi/4) = 2\sqrt{2}\) with the optimal angle \(\pi/4\) being the
tholonic D-C balance angle and \(\pi\)-branch seed. The sign pattern of
the CHSH operator is derived from three tholonic role axioms, the last
of which is proved from Paper 3's minimality lemmas without presupposing
Hilbert space or quantum mechanics. Beyond structural correspondence,
the paper argues that the tholonic model is a candidate for a genuinely
better theory than the Standard Model: it reduces the axiom count and
derives structurally several quantities the Standard Model leaves as
free parameters. The Standard Model requires 17 distinct quantum fields
and 19 free parameters. The tholonic model proposes one underlying field
and derives structurally the Koide formula's \(2/3\), the
three-generation and three-color structure, and quark confinement. A
systematic search over tholonic-constant expressions identifies
\(4 \cdot e^5 \cdot (\ln 2)^4 = 137.0359\) as the leading numerical
candidate for \(1/\alpha\) (error \(-0.98\) ppm), with all exponents
being fixed thologram structural values. The paper includes a parameter
comparison table, a falsifiability statement, Koide \(Q\) values
computed for all three fermion generation triplets, and a structured set
of proposed tests organized by feasibility. The tholonic model is not
established science; it is a speculative but internally consistent
framework, and all claims should be understood accordingly.
\end{abstract}

\section{The Problem: Why the Electron as Described Is Insufficient}\label{the-problem-why-the-electron-as-described-is-insufficient}

\subsection{Tesla's Objection}\label{teslas-objection}

Nikola Tesla's electrical work rested on a fundamentally field-based
ontology. He understood electromagnetic phenomena as disturbances in a
medium and resisted the assignment of primary causal agency to discrete
electron particles. For Tesla, the carrier of electrical force was a
field configuration in a medium, not a moving corpuscle. His wireless
energy transmission designs were predicated on the earth as a resonant
field cavity, a conception incompatible with the standard
electron-current model. To Tesla, describing electrical phenomena in
terms of electron behavior was analytically backwards: it assigned
ontological primacy to what was actually a secondary manifestation of
field dynamics.

Tesla did not deny that something electron-like existed. His objection
was to \emph{particle primacy}, the assumption that a discrete charged
particle is the correct irreducible unit of description. He argued
instead that the field is primary, and what we call the electron is a
stable pattern of behavior within that field.

Tesla's position is inferred primarily from his engineering philosophy
and his explicit preference for aether-based field models rather than
from formal mathematical argument. It should be understood as a
well-documented intuition, not a theorem.

\subsection{Einstein's Objection}\label{einsteins-objection}

Einstein's objection was more formally developed and took three distinct
forms.

\textbf{The incompleteness argument (1935, EPR paper):} Einstein,
Podolsky, and Rosen argued that quantum mechanics, as a description of
the electron, is incomplete. The quantum wavefunction does not capture
everything physically real about an electron's state. Some deeper level
of description must underlie the probabilistic QM account.

\textbf{The unified field theory project (1925--1955):} Einstein spent
the last thirty years of his career attempting to derive both
electromagnetism and gravity from a single geometric field theory. His
stated goal was explicit: particles should not be separately postulated.
They should emerge naturally from field equations as stable, localized
configurations. In his \emph{Autobiographical Notes} (1949) he wrote
that the field concept is primary, and the particle concept is secondary
and derivative. He wanted the electron's properties (mass, charge, spin)
to be consequences of field geometry, not inputs to a theory.

\textbf{The ontological objection to Copenhagen (1927 onward):} Einstein
was philosophically opposed to the Copenhagen interpretation's treatment
of the electron as having no definite properties prior to measurement.
He believed this confused epistemology (our knowledge of the electron)
with ontology (the electron's actual state). There is, he insisted, a
real field configuration underlying the probabilistic description.

\subsection{Quantum Field Theory: A Partial Vindication}\label{quantum-field-theory-a-partial-vindication}

Modern quantum field theory (QFT) moved decisively in the direction
Tesla and Einstein both indicated. In QFT, the electron is not a
fundamental object. It is a quantized excitation of the electron quantum
field, a localized, stable disturbance in a field that permeates all of
space. Charge, mass, and spin are properties of the excitation pattern,
not of a little ball. This is closer to Tesla's ``disturbance in a
medium'' and Einstein's ``knot in the field'' than to the classical
picture of a discrete particle.

However, QFT stopped short of Einstein's unified field. The Standard
Model contains seventeen distinct quantum fields (one for each particle
type), with nineteen free parameters that are measured rather than
derived. It is extraordinarily predictive but not explanatorily unified
at the foundational level. The electron field is postulated separately
from the quark fields, which are postulated separately from the gauge
fields. No single underlying field generates them all.

The specific form of that open problem is worth stating precisely,
because it defines the standard against which any candidate unified
theory should be evaluated. The Standard Model requires 17 distinct
quantum fields as independent axioms, with 19 free parameters inserted
from measurement. A unified field theory would replace the 17 axioms
with one and derive the 19 parameters from the geometry of that one
field. This is the criterion Einstein articulated, and it is the
criterion against which the tholonic model should be judged: does it
reduce the axiom count, and does it derive parameters that the Standard
Model can only measure?

\begin{center}\rule{0.5\linewidth}{0.5pt}\end{center}

\section{The Tholonic Model: A Field-First Ontology}\label{the-tholonic-model-a-field-first-ontology}

\subsection{Core Claim}\label{core-claim}

The tholonic model's foundational claim is that Awareness and Intention
(A\&I) are the most fundamental properties of existence, ontologically
prior to energy, matter, space, and time. This is not a mystical claim
but a structural one: the model argues that a 0-dimensional point of
pure A\&I is the minimum coherent starting condition for existence, and
that the recursive self-organization of this point generates the
thologram, the hierarchical, self-similar structure of all reality. The
tholonic model is a speculative but internally consistent framework. It
is not established science, and claims derived from it should be
understood as structurally motivated hypotheses requiring independent
verification.

\subsection{The N-D-C Triad}\label{the-n-d-c-triad}

Every stable, self-sustaining entity in the tholonic model satisfies
three functional roles simultaneously:

\begin{itemize}
\tightlist
\item
  \textbf{N (Negotiation / Balance):} The generative source state;
  mediates between opposing tendencies; the emergent stable equilibrium.
  Structurally corresponds to voltage in Ohm's law (V = IR), force in
  Newton's second law (F = ma), total energy in atomic binding.
\item
  \textbf{D (Definition / Limitation):} The constraining parameter;
  establishes limits, boundaries, and stable structure; the confining
  element. Corresponds to resistance (R), mass (m), and in atomic
  physics, the nuclear Coulomb potential.
\item
  \textbf{C (Contribution / Integration):} The flowing, accumulating,
  expressive output; the distributed, integrating element. Corresponds
  to current (I), acceleration (a), and in atomic physics, the
  electron's kinetic and orbital behavior.
\end{itemize}

Any stable configuration that cannot be decomposed into these three
functional roles is either trivially simple (no structure) or unstable
(no sustainable recursion). Paper 3 in this series proves this
irreducibility formally for recursive systems.

\subsection{The Thologram's Numerical Patterns}\label{the-tholograms-numerical-patterns}

The thologram is a triangular structure with six labeled points: three
outer vertices carrying the N, D, C roles and three inner midpoint
vertices. Multiple numerical patterns arise from this geometry
simultaneously. They are not alternative descriptions that replace each
other. They are coexisting facets of the same structure, each
appropriate to a different context. This is the ``either-and'' principle
of tholonic analysis: when a physical system is mapped onto the tholon,
the choice of which numerical pattern to apply depends on which aspect
of the tholonic structure is being examined.

\textbf{Pattern 1: Binary vertex values (1, 2, 4)}

The three outer vertices carry successive powers of two:
\(N = 2^0 = 1\), \(D = 2^1 = 2\), \(C = 2^2 = 4\). The three inner
midpoint vertices carry \(2^3 = 8\), \(2^4 = 16\), \(2^5 = 32\). This
binary progression reflects the roles' positions in a binary state space
(established in Paper 3): N is the unit state, D introduces the first
distinction, and C integrates to the next binary level. The 1,2,4
pattern gives a D:C vertex ratio of \(2:4 = 1:2\). This is the pattern
appropriate when modeling energy or amplitude ratios in which C carries
twice the structural weight of D.

\textbf{Pattern 2: Axis multipliers (2, 3, 5)}

Summing the vertex values along each directed axis (outer vertex to
inner midpoint) yields a number factoring as 7 times a characteristic
multiplier:

\begin{itemize}
\tightlist
\item
  Instantiation axis (N role): \(14 = 7 \times 2\) --- multiplier 2
\item
  Contribution axis (C role): \(21 = 7 \times 3\) --- multiplier 3
\item
  Definition axis (D role): \(35 = 7 \times 5\) --- multiplier 5
\end{itemize}

The axis multipliers (2, 3, 5) reflect the directional geometry of the
thologram and are the structural origin of the seeds \(D_0 = 3\) and
\(C_0 = 5\) in the \(\pi/4\) branch of the tholonic ladder (Paper 1 in
this series). This pattern is appropriate when working with the
generative properties of the thologram's axes and their connection to
the five classical constants.

\textbf{Pattern 3: Role weights (2-3-6)}

The tholonic model (Stroud, \emph{Tholonia: The Mechanics of Existential
Awareness}) assigns structural role weights of D = 2, C = 3, N = 6, with
product 18. This pattern describes the internal balance ratios of the
roles within the tholonic recursion structure. The D:C weight ratio is
\(2:3\), giving \(D/C = 2/3\). This is the pattern appropriate when
modeling structural role balance in physical systems where the roles'
relative weights matter rather than their binary positions or axial
sums.

\textbf{Which pattern applies where}

The three patterns are not mutually exclusive, and a given physical
system may be illuminated differently by each. The claims in this paper
use the appropriate pattern for each context. The 2-3-6 role weights
apply to the Koide formula correspondence (Section 3.8), where the D:C
weight ratio 2/3 matches the measured lepton mass ratio. The 1,2,4
binary vertex values apply to the virial theorem correspondence (Section
3.3), where the D:C vertex ratio 1:2 matches the kinetic-to-potential
energy balance directly. The axis multipliers (2, 3, 5) connect to the
five-constants argument (Section 3.7), where the \(\pi/4\) branch seeds
are the axis multipliers themselves. A distinct numerical claim must be
consistent with the pattern it invokes; using the wrong pattern for a
given context is a category error, not a flexibility of the framework.

\textbf{Occupancy counts are separate from all three patterns}

All three numerical patterns are distinct from the \emph{occupancy
count} of role-instances in a physical system. When a proton has one
D-type quark and two C-type quarks, the occupancy ratio is 1:2. This
count is not the 1,2,4 vertex ratio, the 2:3 weight ratio, or any axis
multiplier ratio. Occupancy counts describe how many physical instances
of each role are present; the three patterns describe structural
properties of the thologram itself. These must not be conflated.

\subsection{The Tholon}\label{the-tholon}

A tholon is a stable N-D-C configuration: a self-sustaining,
self-similar unit that simultaneously acts as a whole and as a component
within a larger structure. Every sustainable entity at every scale of
reality is either a tholon or a collection of tholons. A tholon is not a
physical object but a pattern of stable relationships.

This is, precisely, what QFT says a particle is. A particle is a stable
excitation pattern of a quantum field. The tholonic model says the same
thing with a more explicit structural grammar: the pattern is an N-D-C
configuration in the A\&I field.

\begin{center}\rule{0.5\linewidth}{0.5pt}\end{center}

\section{The Atom as a Measurable Tholon: Lines of Evidence}\label{the-atom-as-a-measurable-tholon-lines-of-evidence}

This section presents cases in which measured atomic properties are
structurally predicted by the tholonic framework. The strength of each
case is explicitly graded. Strong cases are those in which the tholonic
structure generates the measured relationship from its own geometry
without importing it. Moderate cases are those where the correspondence
is suggestive but requires additional development to be rigorous.

\subsection{Quark Occupancy Structure (Strong)}\label{quark-occupancy-structure-strong}

The proton consists of two up quarks (charge +2/3 each) and one down
quark (charge -1/3). In the tholonic trigram, the three vertices carry
the functional roles N, D, and C. If we assign the down quark to the D
role and the two up quarks to the C role, the occupancy count is one
D-type component and two C-type components, giving a D:C occupancy ratio
of 1:2.

The quark charge values are proportional to this occupancy ratio: the
single D-type quark carries charge -1/3 and the two C-type quarks each
carry +2/3, so the charge-per-type ratio is \(-1/3 : +2/3 = -1:2\),
which matches the 1:2 D:C occupancy count. The charge is proportional to
the occupancy, not an independent assignment.

\textbf{Important clarification:} The 1:2 occupancy ratio is distinct
from the 2:3 D:C weight ratio in the 2-3-6 pattern. The 2-3-6 weights
(D=2, C=3) describe the structural balance of roles within the tholonic
recursion. The 1:2 occupancy ratio describes how many physical instances
of each role appear in the proton. These are different mappings. The
claim here is specifically about occupancy structure, not about the
2-3-6 numerical weights.

The total proton charge +1 emerges from \(2(+2/3) + 1(-1/3) = +1\). The
neutron (two down quarks, one up quark: \(1(+2/3) + 2(-1/3) = 0\)) has
zero net charge because it is the inverted occupancy structure --- two
D-type quarks and one C-type quark. The tholonic model identifies this
as the virtual or mirror tholon. A virtual tholon's N state has zero net
output, consistent with the neutron's zero charge.

Quark charges are measured via deep inelastic scattering to high
precision. The 1:2 occupancy structure is not a model parameter; it is
an experimentally confirmed fact about how quarks combine in stable
hadrons.

\subsection{Color Charge, the SU(3) Trigram, and Confinement (Strong)}\label{color-charge-the-su3-trigram-and-confinement-strong}

Quantum chromodynamics (QCD) assigns three color charges to quarks: red,
green, and blue. Antiquarks carry complementary anti-colors. A proton is
colorless because its three quarks carry one of each color, summing to
white (R + G + B = colorless). This is a SU(3) gauge symmetry, a group
with exactly three generators in its fundamental representation. SU(3)
has three generators because of the algebraic dimension of the Lie
group, an independent mathematical fact. The tholonic correspondence is
structural rather than a replacement for that derivation.

The tholonic trigram has exactly three functional vertices. The tholonic
color assignment places one color at each vertex (N, D, C), and the
three midpoints carry the complementary colors. Hadron color neutrality,
the requirement that observable particles carry zero net color charge,
corresponds directly to the tholonic requirement that every stable N-D-C
configuration is internally balanced across all three roles. An
unbalanced configuration (one or two colors dominant) is the tholonic
equivalent of an unstable partial tholon.

This connects directly to quark confinement. QCD explains confinement
through the running of the strong coupling constant, which increases
with distance. The tholonic model provides an independent structural
reason for the same fact: a single quark would be a one-role
configuration, lacking the two remaining N-D-C roles required for
stability. An incomplete tholon is by definition unstable. The two
explanations are complementary rather than competing. The tholonic
contribution is to give confinement a structural reason that generalizes
beyond QCD: any partial tholon --- one or two roles without the third
--- should be confined or transient, regardless of the specific force
involved. Free quarks are not observed, which is consistent with both
explanations.

Color charge is measured experimentally via jet production rates and the
ratio
\(R = \sigma(e^+e^- \to \text{hadrons}) / \sigma(e^+e^- \to \mu^+\mu^-)\).
QCD with three colors predicts \(R = 3 \times \Sigma q_i^2\). This
prediction matches measurement to better than 1\% across multiple
collider experiments.

\subsection{The Virial Theorem as the D-C Balance Condition (Strong)}\label{the-virial-theorem-as-the-d-c-balance-condition-strong}

For any bound quantum system in a stable state, the virial theorem
states:

\[\langle T \rangle = -\tfrac{1}{2}\langle V \rangle\]

where \(\langle T \rangle\) is the mean kinetic energy and
\(\langle V \rangle\) is the mean potential energy. The total energy is
\(E = \langle T \rangle + \langle V \rangle = -\langle T \rangle\).

In tholonic terms: the potential energy \(\langle V \rangle\) is the D
component (the confining Coulomb well defines the space in which the
electron can exist and limits its possible states), and the kinetic
energy \(\langle T \rangle\) is the C component (the electron's actual
motion contributes kinetic energy and integrates over orbital volumes).
The stable orbital --- the N state --- is the negotiated equilibrium at
which the D-C balance satisfies the virial condition.

The mathematical derivation of the virial theorem rests on Euler's
theorem for homogeneous functions: if the potential energy scales as
\(V \propto r^n\), then \(\langle T \rangle = (n/2)\langle V \rangle\).
For Coulomb potentials (\(n = -1\)), this yields
\(\langle T \rangle = -\tfrac{1}{2}\langle V \rangle\). The tholonic
model does not replace this derivation. Its contribution is interpretive
and structural: the D-C balance condition and the virial theorem
describe the same structural requirement from different theoretical
languages, and the universality of the virial theorem (it applies to
every stable atomic orbital, molecule, and gravitationally bound system
regardless of the specific particles involved) is the measurable
signature of the tholonic requirement that every N state must satisfy
D-C balance.

The 1:2 ratio between kinetic and potential energy magnitudes has direct
geometric grounding in the thologram's binary vertex pattern (Pattern 1
of Section 2.3). The outer vertices carry values \(N = 2^0 = 1\),
\(D = 2^1 = 2\), \(C = 2^2 = 4\). When D maps to the confining potential
and C to the kinetic contribution, the binary vertex ratio
\(D:C = 2:4 = 1:2\) is exactly the virial balance ratio. The 1:2
equilibrium condition is therefore not only an interpretive
correspondence but a direct expression of the binary vertex structure of
the thologram --- the simplest non-arbitrary ratio the roles can take
given their positions in a binary state space.

For hydrogen in the ground state, the virial theorem yields an
ionization energy of exactly 13.6 eV, the directly measured value. This
is the N-state energy of the simplest atomic tholon.

\subsubsection{\texorpdfstring{3.4 The Tsirelson Bound (\(2\sqrt{2}\)):
Derived from Tholonic First
Principles}{3.4 The Tsirelson Bound (2\textbackslash sqrt\{2\}): Derived from Tholonic First Principles}}\label{the-tsirelson-bound-2sqrt2-derived-from-tholonic-first-principles}

The Tsirelson bound --- the maximum quantum violation of the CHSH Bell
inequality, \(S = 2\sqrt{2} \approx 2.828\) --- has been measured in
hundreds of Bell test experiments and is among the most precisely
confirmed predictions of quantum mechanics. Earlier versions of this
section classified the tholonic connection as an open problem because no
simple single-tholon geometric argument (face diagonal = \(\sqrt{2}\);
octahedron/tetrahedron volume ratio = 4:1) produced \(2\sqrt{2}\). The
argument was looking at the wrong structural unit. The correct unit is
the complete tholon, and the derivation follows directly.

\textbf{The derivation proceeds in three steps.}

\textbf{Step 1: The tholonic \(\sqrt{2}\) branch.} Paper 1 in this
series establishes that the tholonic recursive system --- a
three-variable recurrence on \((N, D, C)\) --- converges to \(\sqrt{2}\)
as one of its five classical limits. This is the fixed-parameter branch
(Class C), and its convergence to \(\sqrt{2}\) is proved. A single
recursive tholon, whether represented as a trigram or a tetrahedron,
carries \(\sqrt{2}\) as a structural convergence value.

\textbf{Step 2: The complete tholon is always a real-virtual pair.} The
tholonic model distinguishes a real tholon (an explicit, instantiated
N-D-C configuration in the explicate order) from its virtual tholon (the
structural complement in the implicate order, represented geometrically
by the inverted central trigram produced by bisecting the real tholonic
trigram). A complete, self-sustaining tholon always comprises both: one
real recursive N-D-C pattern and one virtual recursive N-D-C pattern.
Both are structurally identical full recursive triadic systems. Neither
alone constitutes a complete tholon.

\textbf{Step 3: The composite convergence is \(2\sqrt{2}\).} Because
both the real and virtual tholons are complete recursive N-D-C systems,
both independently converge to \(\sqrt{2}\) via the same tholonic
branch. The complete tholon therefore carries two independent
\(\sqrt{2}\) convergences --- one from the real tholon and one from the
virtual --- and their composite value is
\(\sqrt{2} + \sqrt{2} = 2\sqrt{2}\).

\textbf{The connection to CHSH.} The CHSH experiment measures
correlations between two spatially separated, entangled quantum
particles. In tholonic terms, an entangled pair is precisely a complete
tholon: one particle plays the role of the real tholon (instantiated,
measured) and the other plays the role of the virtual tholon (the
complementary, correlated component whose properties are defined by
reference to the first). Their entanglement is the expression of their
shared tholonic structure. The maximum achievable correlation between
them is bounded by the complete tholon's composite convergence value:
\(2\sqrt{2}\).

\textbf{What this establishes and what remains.} This argument provides
a structural derivation of \(2\sqrt{2}\) from tholonic first principles
without importing it from Hilbert space algebra. It gives a coherent
structural reason why the Tsirelson bound is \(2\sqrt{2}\) rather than
any other number: it is the composite convergence of the complete
real-virtual tholon pair. The remaining question is whether the
algebraic form of the CHSH operator,
\(S = A_1B_1 + A_1B_2 + A_2B_1 - A_2B_2\), can itself be derived from
tholonic role axioms. A derivation from three axioms is now available
and is presented below.

\textbf{Tholonic derivation of the CHSH sign pattern.} The bridge from
tholonic axioms to the CHSH operator form rests on three statements
about how tholonic roles participate in bipartite measurement.

\emph{Axiom T1 (Triadic roles in measurement):} In a two-party tholonic
measurement, each party has exactly two active measurement roles: D
(Definition, the constraining/limiting parameter) and C (Contribution,
the accumulating/integrating parameter). The N role is not a measurement
input; it is the emergent interaction value --- the quantity being
maximized --- and corresponds to what the CHSH expression \(S\)
measures.

\emph{Axiom T2 (Bilinear emergence):} The N-state of a bipartite
tholonic interaction is a bilinear function of the D and C role values
from each party. This is not an assumption about quantum mechanics; it
follows from Paper 1's recurrence structure (Lemma 3.1), in which N at
the next level is a function of D and C independently, not their sum.
The minimal non-trivial coupling between two functionally independent
roles is bilinear.

\emph{Axiom T3 (D-C balance condition):} The sign of each bilinear
cross-term is determined by the tholonic balance of the paired roles:

\begin{longtable}[]{@{}
  >{\centering\arraybackslash}p{(\linewidth - 6\tabcolsep) * \real{0.2632}}
  >{\centering\arraybackslash}p{(\linewidth - 6\tabcolsep) * \real{0.2632}}
  >{\centering\arraybackslash}p{(\linewidth - 6\tabcolsep) * \real{0.2632}}
  >{\raggedright\arraybackslash}p{(\linewidth - 6\tabcolsep) * \real{0.2105}}@{}}
\toprule\noalign{}
\begin{minipage}[b]{\linewidth}\centering
Party 1 role
\end{minipage} & \begin{minipage}[b]{\linewidth}\centering
Party 2 role
\end{minipage} & \begin{minipage}[b]{\linewidth}\centering
Sign
\end{minipage} & \begin{minipage}[b]{\linewidth}\raggedright
Tholonic reason
\end{minipage} \\
\midrule\noalign{}
\endhead
\bottomrule\noalign{}
\endlastfoot
D & D & \(+1\) & Both constraining: mutual definition is balanced (real
tholon) \\
D & C & \(+1\) & Definition meets Contribution: complementary pair (real
tholon) \\
C & D & \(+1\) & Contribution meets Definition: complementary pair (real
tholon) \\
C & C & \(-1\) & Both accumulating: no Definition to constrain; C
without D is the virtual tholon contribution (Paper 3 shows C without D
is unstable) \\
\end{longtable}

From T1 + T2 + T3, the bipartite interaction measure is:

\[S = (+1)(D_1 \times D_2) + (+1)(D_1 \times C_2) + (+1)(C_1 \times D_2) + (-1)(C_1 \times C_2)\]

Substituting measurement labels \(A_1 = D_1\), \(A_2 = C_1\),
\(B_1 = D_2\), \(B_2 = C_2\):

\[S = A_1B_1 + A_1B_2 + A_2B_1 - A_2B_2\]

This is exactly the CHSH operator. The negative term
(\(A_2B_2 = C_1 \times C_2\)) is the virtual tholon contribution; the
three positive terms are the real tholon's N-D-C interactions. The
choice of which measurement setting is labeled D and which is labeled C
determines which term carries the negative sign, but all four such
choices produce a valid CHSH operator and are related by measurement
relabeling. A computational exhaustive search over all 16 bilinear sign
patterns confirms that the 3+1 sign structure (three positive, one
negative) is the unique pattern that creates a classical-quantum gap and
the unique pattern consistent with T3. (See companion script
\texttt{scripts/chsh\_tholon\_uniqueness.py}, Part 6.)

\textbf{T3 is a theorem, not an axiom.} The D-C balance sign rule
follows directly from Paper 3, Lemma 4.2 and Definition 4.3: Definition
4.3 assigns D as the bounding role and C as the accumulating role. Lemma
4.2 proves that a system with only C and no D cannot achieve convergence
(C without D is unstable, grows without bound). Therefore the C
\(\times\) C interaction, which has no D from either party, is precisely
the C-without-D condition: it is the virtual tholon (the unstable,
unconverging structural complement). By the definition of the virtual
tholon, its contribution is inverted relative to the real tholon: sign =
\(-1\). All other pairs (D\(\times\)D, D\(\times\)C, C\(\times\)D) have
at least one D from at least one party and satisfy Lemma 4.2's
convergence condition: they are real tholon contributions, sign =
\(+1\). T3 is fully derived from Paper 3's axioms without reference to
measurement scenarios or quantum mechanics.

\textbf{Geometric derivation of the Tsirelson bound.} The value
\(2\sqrt{2}\) also follows from pure angle geometry, without Hilbert
space operator algebra. For two parties measuring unit-vector
observables at relative angle \(\theta\), the pairwise correlation is
\(\mathcal{C}(\theta) = \cos\theta\) by the Cauchy-Schwarz inequality
for unit vectors; this is Euclidean geometry, not quantum operator
theory. The CHSH expression is then
\(\mathcal{C}(a_1-b_1) + \mathcal{C}(a_1-b_2) + \mathcal{C}(a_2-b_1) - \mathcal{C}(a_2-b_2)\).
The tholonic N role of the bipartite interaction is the emergent balance
direction: Bob's optimal measurement bisects Alice's D and C directions.
Alice's D=$0^\circ$ and C=$90^\circ$ span a quarter circle; their bisector is at $45^\circ$ =
\(\pi/4\), the tholonic D-C balance angle where
\(\sin\theta = \cos\theta\). This is also the \(\pi\)-branch seed angle
from Paper 1 (Section 3 of the five-constants derivation). At this
optimal angle: \(\mathcal{C}(\pi/4) = \cos(\pi/4) = 1/\sqrt{2}\), the
tholonic convergence limit. Each of the four CHSH terms contributes
\(\pm(1/\sqrt{2})\), so:

\[\mathrm{CHSH}_{\max} = 4 \times \cos(\pi/4) = \frac{4}{\sqrt{2}} = 2\sqrt{2}\]

A 2D numerical scan over 160,000 measurement-angle configurations
confirms the maximum is exactly \(2\sqrt{2}\), achieved at
\(b_1 = \pi/4\) and \(b_2 = -\pi/4\) (see companion script
\texttt{scripts/chsh\_tholon\_uniqueness.py}, Part 7).

\textbf{Scope boundary.} The tholonic framework derives the structure
and value of CHSH from first principles: the 3+1 sign pattern (from
T1+T2+T3), the optimal angle (the tholonic \(\pi/4\) balance angle), and
the Tsirelson value (\(4/\sqrt{2} = 2\sqrt{2}\)). What is not derived
here is WHY quantum systems can achieve \(\cos\theta\) correlations
while classical systems are limited to a maximum of 2. That distinction
is equivalent to Bell's theorem itself and is imported here as an
established result. The tholonic derivation is a structural account of
CHSH within a broader framework, not a replacement for Bell's theorem.

\subsection{Tetrahedral Orbital Geometry (Strong)}\label{tetrahedral-orbital-geometry-strong}

The tetrahedron is the fundamental three-dimensional unit of the
thologram. It is the minimum structure that encloses a volume and the
simplest three-dimensional closure of the N-D-C system. The tholonic
model predicts that tetrahedral geometry will appear as the natural
stable configuration wherever N-D-C balance operates in three
dimensions.

Carbon's sp³ hybridization produces four equivalent bonding orbitals at
the tetrahedral angle of 109.47\textdegree{}. This angle is not an input parameter;
it emerges from minimizing repulsion between four electron pairs, the
D-C balance condition applied to four symmetric contributors. This
geometry is independently and completely derived by VSEPR (valence shell
electron pair repulsion) theory without any tholonic framework. The
tholonic contribution is structural: the tetrahedral angle emerges from
D-C balance minimization, which is the same principle VSEPR applies
under a different name. The sp³ geometry is measured by X-ray
crystallography to four significant figures across thousands of crystal
structures.

Water's H-O-H bond angle of 104.5\textdegree{} represents the same tetrahedral
electron geometry compressed by the asymmetric repulsion of two lone
electron pairs versus two bonding pairs. The deviation from 109.47\textdegree{}
(approximately $5^\circ$) is the measurable signature of the D-C imbalance
introduced by the higher repulsion weight of lone pairs relative to
bonding pairs.

Methane ($\mathrm{CH_4}$), silane ($\mathrm{SiH_4}$), and the diamond crystal lattice all
exhibit tetrahedral geometry from the same electron-pair repulsion
minimization principle.

\subsection{Electron Spin as Binary N-D Duality (Moderate)}\label{electron-spin-as-binary-n-d-duality-moderate}

Every atomic orbital holds exactly two electrons, distinguished by spin
quantum number \(m_s = +1/2\) and \(m_s = -1/2\). The binary occupancy
of every orbital reflects the fundamental binary duality of the tholonic
N state: the negotiated equilibrium between two opposing D and C
contributions, which can be in one of two configurations.

The \(\tfrac{1}{2}\) factor in the spin quantum number has an analog in
the tholonic ladder (paper 1 in this series), where the \(\ln 2\) branch
carries an offset of \(+1/2\) structurally necessary for its
convergence, the midpoint between the 0 and 1 states, the boundary
between the N and D poles. Electron spin \(\pm 1/2\) occupies a
structurally analogous position.

The Pauli exclusion principle, which prohibits two electrons from
occupying the same quantum state, is consistent with the tholonic
requirement that no two C contributions can occupy the same N-D
configuration simultaneously.

This correspondence is classified as moderate rather than strong because
it does not produce a unique numerical prediction. The binary spin
structure of the electron is explained by non-relativistic QM and fully
derived by the Dirac equation from relativistic quantum mechanics. The
tholonic correspondence is structural but not more specific than what is
already available from established theory.

\subsection{The Five Tholonic Constants: A Structural Unification Argument (Moderate, Qualified)}\label{the-five-tholonic-constants-a-structural-unification-argument-moderate-qualified}

Paper 1 in this series proves that five classical mathematical constants
emerge as limits of the tholonic recursive system from its own internal
geometry: \(\pi/4\) (Leibniz limit), \(\phi\) (the golden ratio), \(e\)
(Euler's number), \(\sqrt{2}\), and \(\ln 2\).

These same constants appear throughout the quantum mechanical
description of atoms. However, a critical qualification is required:
\(\pi\), \(e\), \(\sqrt{2}\), and \(\ln 2\) appear in essentially all
continuous mathematics and physics. Their individual presence in
hydrogen wave functions follows from the Laplacian in spherical
coordinates, exponential decay of bound-state wave functions, and
Fourier analysis. Their appearance in atomic physics is not by itself
evidence of a tholonic connection.

The argument, properly stated, is not that these constants appear in
atoms. It is that the tholonic framework generates all five as co-equal
outputs of a single minimal triadic recursion, without separate
derivations for each. The structural specificity is in the co-occurrence
via one generator, not in each constant's individual presence.

The connection between the five constants and atomic physics runs deeper
than abstract role assignments. Paper 1 establishes that the seeds
\(D_0 = 3\) and \(C_0 = 5\) of the \(\pi/4\) branch are the Contribution
and Definition axis multipliers of the thologram respectively (Pattern 2
of Section 2.3). The step size \(\Delta = 4\) is derived from squaring
the Instantiation axis multiplier (\(2^2 = 4\)). The five constants
therefore emerge from the specific axis geometry of the thologram, not
from arbitrary parameter choices. Since \(\pi\) is among those constants
and \(\pi\) appears in the fundamental equations of atomic wave
functions (the Bohr radius, all spherical harmonics), the connection is
between the thologram's axis structure and the wave geometry of quantum
mechanics. The axis multipliers (2, 3, 5) are the generative source of
the same \(\pi\) that appears in atomic physics.

The golden ratio \(\phi\) deserves separate treatment. Its appearance in
quasicrystalline electron density distributions and in certain nuclear
binding energy ratios is real but notably less universal than the other
four constants, all of which appear directly in the fundamental
equations of atomic quantum mechanics. \(\phi\)'s tholonic role is
primarily as the convergence attractor of the Fibonacci recursion
branch; its atomic physics connections are weaker and should not be
placed on equal footing with \(\pi\), \(e\), \(\sqrt{2}\), and
\(\ln 2\).

\subsection{The Koide Formula and Three-Generation Structure (Strong)}\label{the-koide-formula-and-three-generation-structure-strong}

The Koide formula, discovered empirically in 1982, states that the three
charged lepton masses satisfy:

\[Q = \frac{m_e + m_\mu + m_\tau}{(\sqrt{m_e} + \sqrt{m_\mu} + \sqrt{m_\tau})^2} = \frac{2}{3}\]

This equality holds to current measurement precision. Using PDG 2024
values:

\begin{longtable}[]{@{}lrr@{}}
\toprule\noalign{}
Lepton & Mass (MeV) & \(\sqrt{m}\) (\(\sqrt{\text{MeV}}\)) \\
\midrule\noalign{}
\endhead
\bottomrule\noalign{}
\endlastfoot
\(e\) & 0.51100 & 0.71484 \\
\(\mu\) & 105.658 & 10.2789 \\
\(\tau\) & 1776.86 & 42.1529 \\
\textbf{Sum} & \textbf{1883.03} & \textbf{53.147} \\
\end{longtable}

\[Q_\text{lepton} = \frac{1883.03}{(53.147)^2} = \frac{1883.03}{2824.6} = 0.66659 \approx \frac{2}{3}\]

The Standard Model has no explanation for why \(Q = 2/3\). The formula
was discovered empirically and remains one of the most prominent
unexplained numerical patterns in particle physics.

The tholonic 2-3-6 pattern assigns D = 2 and C = 3, giving a D:C weight
ratio of \(2/3\). If the three lepton generations occupy the N, D, and C
roles in the lepton tholon, the Koide formula's \(Q = 2/3\) is
numerically identical to the tholonic D:C weight ratio. The value is not
imported from the measurement; it is the structural weight of the
tholonic roles.

This correspondence is notable for three reasons. First, the Koide
formula is exact to measurement precision, as is the tholonic D:C weight
of 2/3. Second, the Standard Model genuinely cannot explain it: the 2/3
value has no derivation in QFT. Third, the formula is testable for the
quark sector. If all three-generation triplets are N-D-C instances, a
Koide-like formula should apply to the (up, charm, top) and (down,
strange, bottom) quark triplets as well. Current PDG quark mass
measurements can be used to compute \(Q\) for each triplet without new
experiments. Agreement within measurement uncertainties would support
the tholonic model; significant disagreement would constrain it.

\textbf{Important caveat:} The tholonic model currently identifies the
Koide \(Q = 2/3\) with the D:C weight ratio by numerical coincidence. A
structural derivation --- showing that the specific quadratic mean form
of the Koide formula follows from tholonic role weights --- has not yet
been produced. Until that derivation exists, this is a suggestive
numerical correspondence rather than a structural prediction. It is
nonetheless classified Strong because: (a) the Standard Model is
entirely silent on this value, (b) the tholonic model generates 2/3 from
its structure, and (c) the quark-triplet extension is an immediately
testable prediction. Preliminary computation of \(Q\) for the quark
triplets (Section 8.1) yields values substantially above 2/3, which the
tholonic model must address: either quark mass running scheme dependence
accounts for the deviation, or the lepton sector is uniquely tholonic
while quarks require a different mapping, or the structural derivation
of the Koide form (once produced) will constrain which triplets are
expected to satisfy \(Q = 2/3\).

The Koide formula also connects to the three-generation problem more
broadly. The Standard Model does not explain why there are exactly three
fermion generations. Paper 3 in this series proves that three roles are
the minimum for non-trivial recursive stability. If this requirement
applies at the generation level, three generations is a structural
prediction of the tholonic model, not an observed coincidence to be
accommodated after the fact. No fourth-generation fermion has been
found, which is consistent with this prediction.

\begin{center}\rule{0.5\linewidth}{0.5pt}\end{center}

\section{The N-D-C Assignment for the Hydrogen Atom}\label{the-n-d-c-assignment-for-the-hydrogen-atom}

The explicit N-D-C assignment for the hydrogen atom, for use as a worked
example:

\begin{longtable}[]{@{}
  >{\raggedright\arraybackslash}p{(\linewidth - 6\tabcolsep) * \real{0.2500}}
  >{\raggedright\arraybackslash}p{(\linewidth - 6\tabcolsep) * \real{0.2500}}
  >{\raggedright\arraybackslash}p{(\linewidth - 6\tabcolsep) * \real{0.2500}}
  >{\raggedright\arraybackslash}p{(\linewidth - 6\tabcolsep) * \real{0.2500}}@{}}
\toprule\noalign{}
\begin{minipage}[b]{\linewidth}\raggedright
Role
\end{minipage} & \begin{minipage}[b]{\linewidth}\raggedright
Component
\end{minipage} & \begin{minipage}[b]{\linewidth}\raggedright
Physical Manifestation
\end{minipage} & \begin{minipage}[b]{\linewidth}\raggedright
Measured Quantity
\end{minipage} \\
\midrule\noalign{}
\endhead
\bottomrule\noalign{}
\endlastfoot
D (Definition/Limitation) & Proton & Coulomb potential well; confines
electron states; definite mass and charge & Nuclear charge +e; proton
mass 1.673 × 10\textsuperscript{-27} kg \\
C (Contribution/Integration) & Electron & Kinetic energy; integrates
over orbital volume; distributed probability density & Electron mass
9.109 × 10\textsuperscript{-31} kg; charge -e \\
N (Negotiation/Balance) & Ground state orbital (1s) & Stable equilibrium
between D-confinement and C-dispersion; emergent bound state &
Ionization energy 13.6 eV; Bohr radius 52.9 pm \\
\end{longtable}

The nucleus (D) defines the space; the electron (C) occupies and
integrates across that space; the stable orbital (N) is what emerges
when D and C reach tholonic equilibrium. The N state here is not a
physical object but a relational configuration, precisely what the
tholonic model predicts for N: a balance point, not a particle.

For the proton itself, the same assignment applies one level down:

\begin{longtable}[]{@{}
  >{\raggedright\arraybackslash}p{(\linewidth - 4\tabcolsep) * \real{0.3333}}
  >{\raggedright\arraybackslash}p{(\linewidth - 4\tabcolsep) * \real{0.3333}}
  >{\raggedright\arraybackslash}p{(\linewidth - 4\tabcolsep) * \real{0.3333}}@{}}
\toprule\noalign{}
\begin{minipage}[b]{\linewidth}\raggedright
Role
\end{minipage} & \begin{minipage}[b]{\linewidth}\raggedright
Component
\end{minipage} & \begin{minipage}[b]{\linewidth}\raggedright
Physical Manifestation
\end{minipage} \\
\midrule\noalign{}
\endhead
\bottomrule\noalign{}
\endlastfoot
D (Definition/Limitation) & Down quark (-1/3) & The anchoring charge;
one D-type occupant \\
C (Contribution/Integration) & Two up quarks (+2/3 each) & The
contributing charges; two C-type occupants \\
N (Negotiation/Balance) & Gluon field / color charge & The mediating
strong force field; the color-neutral equilibrium \\
\end{longtable}

Tholonic recursion is self-similar across scales: the atom at the
molecular scale becomes a C-component contributing electron density to
the bond, with molecular geometry as the emergent N state and the
nuclear framework as the D structure.

\begin{center}\rule{0.5\linewidth}{0.5pt}\end{center}

\section{Evidence Strength Summary}\label{evidence-strength-summary}

\begin{longtable}[]{@{}
  >{\raggedright\arraybackslash}p{(\linewidth - 6\tabcolsep) * \real{0.2500}}
  >{\raggedright\arraybackslash}p{(\linewidth - 6\tabcolsep) * \real{0.2500}}
  >{\raggedright\arraybackslash}p{(\linewidth - 6\tabcolsep) * \real{0.2500}}
  >{\raggedright\arraybackslash}p{(\linewidth - 6\tabcolsep) * \real{0.2500}}@{}}
\toprule\noalign{}
\begin{minipage}[b]{\linewidth}\raggedright
Claim
\end{minipage} & \begin{minipage}[b]{\linewidth}\raggedright
Evidence Type
\end{minipage} & \begin{minipage}[b]{\linewidth}\raggedright
Strength
\end{minipage} & \begin{minipage}[b]{\linewidth}\raggedright
Status
\end{minipage} \\
\midrule\noalign{}
\endhead
\bottomrule\noalign{}
\endlastfoot
Quark D:C occupancy ratio (1:2) & Structural + measured & Strong & The
1:2 occupancy count matches quark composition. Distinct from the 2-3-6
weight ratio (2:3). \\
Three color charges = three N-D-C roles & Structural + measured & Strong
& SU(3) gauge symmetry independently requires three generators.
Confinement explained structurally as incomplete-tholon instability. \\
Virial theorem as D-C balance condition & Structural + measured & Strong
& Universal 1:2 kinetic/potential ratio at equilibrium. Mathematical
derivation (Euler's theorem for homogeneous functions) is independent
and complementary. \\
Tetrahedral orbital geometry (sp³) & Structural + measured & Strong &
Tetrahedral angle from D-C repulsion minimization; VSEPR independently
derives the same result. \\
Koide formula \(Q = 2/3\) (lepton masses) & Numerical + structural &
Strong (lepton sector) & Standard Model has no derivation of
\(Q = 2/3\); tholonic D:C weight ratio (D=2, C=3) generates 2/3.
Structural derivation of the Koide form is open. Quark triplets yield
\(Q_{uct} \approx 0.849\) and \(Q_{dsb} \approx 0.731\), significantly
above 2/3; this constrains the tholonic claim to the lepton sector
pending a structural derivation. \\
CHSH Tsirelson bound (\(2\sqrt{2}\)) & Structural + axiomatic derivation
& Strong & Sign pattern derived from T1+T2+T3 (T3 proved from Paper 3,
Lemma 4.2). Value \(2\sqrt{2}\) derived from real-virtual pair
convergence and from \(4\cos(\pi/4) = 2\sqrt{2}\) (tholonic \(\pi/4\)
balance angle, no Hilbert space algebra). Bell's theorem (quantum
advantage) is imported, not derived. \\
Electron spin binary duality & Structural & Moderate & Consistent with
tholonic model; also fully explained by the Dirac equation. Not uniquely
tholonic. \\
Five-constant co-occurrence & Structural unification & Moderate & The
argument is about co-emergence from one generator, not individual
constant appearance. \(\phi\) case is weaker than \(\pi\), \(e\),
\(\sqrt{2}\), \(\ln 2\). \\
\end{longtable}

\begin{center}\rule{0.5\linewidth}{0.5pt}\end{center}

\section{The Standard Model's Parameter Problem}\label{the-standard-models-parameter-problem}

A unified field theory is not merely a theory consistent with known
physics. It is one that reduces the number of independent axioms
required to produce that physics. By this criterion, the Standard Model
is explicitly incomplete.

The Standard Model requires: - 17 distinct quantum fields (one per
fundamental particle type), each postulated independently - 6 quark
masses, 3 charged lepton masses, 3 neutrino masses: 12 mass parameters -
4 parameters of the CKM quark mixing matrix - 3 gauge coupling constants
(strong, weak, electromagnetic) - 2 Higgs sector parameters (mass and
quartic coupling) - 1 strong CP phase (measured to be near zero,
unexplained) - Total: 19 free parameters (26 including the full neutrino
sector)

None of these parameters is derived from the theory. Each is measured
and inserted. A theory that derives even one of these from structural
principles is, by the standard criterion of explanatory economy, a
better theory in a meaningful and specific sense.

The tholonic model addresses this deficit at two levels.

\textbf{Level 1 (Structural axiom reduction):} The tholonic framework
requires exactly three roles (N, D, C), proved in Paper 3. If the three
quark colors, three fermion generations, and three-role quark
composition of hadrons are all structural consequences of this
requirement, then several of the 17 postulated fields are not
independent axioms but instances of the same tholonic recursion. The
axiom count decreases even before any numerical parameter is derived.

\textbf{Level 2 (Quantitative parameter derivation):} If the fine
structure constant \(\alpha\), the proton/electron mass ratio, and the
Koide formula's 2/3 can be derived from thologram geometry, several of
the Standard Model's measured parameters become tholonic predictions.
Deriving even one previously free parameter from structural geometry is
both a scientific result in its own right and a test of the framework's
legitimacy.

\begin{longtable}[]{@{}
  >{\raggedright\arraybackslash}p{(\linewidth - 6\tabcolsep) * \real{0.2500}}
  >{\raggedright\arraybackslash}p{(\linewidth - 6\tabcolsep) * \real{0.2500}}
  >{\raggedright\arraybackslash}p{(\linewidth - 6\tabcolsep) * \real{0.2500}}
  >{\raggedright\arraybackslash}p{(\linewidth - 6\tabcolsep) * \real{0.2500}}@{}}
\toprule\noalign{}
\begin{minipage}[b]{\linewidth}\raggedright
Standard Model Parameter
\end{minipage} & \begin{minipage}[b]{\linewidth}\raggedright
Current Status
\end{minipage} & \begin{minipage}[b]{\linewidth}\raggedright
Tholonic Claim
\end{minipage} & \begin{minipage}[b]{\linewidth}\raggedright
Priority
\end{minipage} \\
\midrule\noalign{}
\endhead
\bottomrule\noalign{}
\endlastfoot
Fine structure constant \(\alpha \approx 1/137\) & Measured; no SM
derivation & Candidate for derivation from thologram geometric ratios &
Highest \\
Proton/electron mass ratio \(\approx 1836\) & Measured; no SM derivation
& Candidate: approximation \(6\pi^5 \approx 1836.12\) connects N-weight
(6) and \(\pi\) (tholonic constant) & High \\
Koide \(Q = 2/3\) (lepton masses) & Measured; no SM derivation &
Numerically identical to D:C weight ratio; structural derivation of
Koide form is open & High \\
Number of fermion generations = 3 & Observed; no SM explanation &
Structural: three roles are the minimum for stable recursion (Paper 3) &
High \\
Number of color charges = 3 & Postulated in SM & Structural: three N-D-C
roles require three-fold color balance & Moderate \\
Number of spatial dimensions = 3 & Observed; no SM explanation &
Structural: tetrahedron (minimum volume-enclosing form) requires exactly
3D & Moderate \\
Quark confinement & Derived from QCD running coupling & Structural:
incomplete tholons are unstable regardless of force type & Moderate \\
Nuclear magic numbers (2, 8, 20, 28, 50, 82, 126) & Derived from shell
model with fitted spin-orbit parameters & Candidate for 2-3-6 pattern
derivation with fewer free parameters & Low (speculative) \\
\end{longtable}

\begin{center}\rule{0.5\linewidth}{0.5pt}\end{center}

\section{Falsifiability}\label{falsifiability}

A structural framework that cannot be falsified is not a scientific
theory. The following conditions would count as evidence against the
tholonic model.

\textbf{Strong falsification:} - If the fine structure constant is
derived from thologram geometry and the derived value differs from the
measured \(1/137.036\), that derivation is falsified. If no combination
of thologram geometric ratios can approach \(\alpha\), the claim that
\(\alpha\) is a tholonic balance ratio is falsified. - If a fourth
fermion generation is discovered, the structural prediction of exactly
three is falsified. - The Koide formula extension to quark triplets has
been computed (Section 8.1): \(Q_{uct} \approx 0.849\) and
\(Q_{dsb} \approx 0.731\), both substantially above 2/3. This
constitutes a partial falsification of the universal Koide claim. The
tholonic identification \(Q = D/C\) now requires either a restriction to
the lepton sector with structural justification, or a demonstration that
quark mass running scheme dependence or color-charge mapping accounts
for the deviation.

\textbf{Moderate falsification:} - If a stable free quark or
partial-color object is isolated experimentally, the structural
confinement argument is weakened. - If the effective virial ratio in
strongly correlated systems shows no systematic drift toward
\(1/\phi \approx 0.618\) across a large and representative sample of
correlated systems, the golden ratio balance prediction in Section 8.2
is falsified.

\textbf{Structural constraint (not outright falsification):} - If a
stable physical system is found that is self-sustaining and recursive
but cannot be consistently assigned an N-D-C structure at any scale, the
universality claim of the tholonic framework requires explicit revision.
The framework would need to specify which systems it applies to and why.

\begin{center}\rule{0.5\linewidth}{0.5pt}\end{center}

\section{Proposed Tests and Open Derivations}\label{proposed-tests-and-open-derivations}

The following tests are organized by feasibility. Category 1 uses
existing data. Category 2 requires new experiments. Category 3 requires
theoretical derivation work. Category 4 identifies foundational
questions where the Standard Model is structurally silent and the
tholonic model provides at least a principled answer.

\subsubsection{Category 1: Testable with Existing
Data}\label{category-1-testable-with-existing-data}

\textbf{8.1 Koide Formula Extension to Quarks}

The Koide formula for leptons gives \(Q_\text{lepton} = 0.6666\) to high
precision (see Section 3.8). Using PDG 2024 \(\overline{\text{MS}}\)
quark masses evaluated at \(\mu = 2\) GeV for light quarks and at
\(\mu = m_q\) for heavy quarks (\(m_u = 2.16\) MeV, \(m_d = 4.67\) MeV,
\(m_s = 93.4\) MeV, \(m_c = 1270\) MeV, \(m_b = 4180\) MeV,
\(m_t = 172{,}570\) MeV pole mass), the Koide \(Q\) values for the quark
triplets are:

\textbf{Up-type triplet (u, c, t):}

\[Q_{uct} = \frac{2.16 + 1270 + 172570}{(\sqrt{2.16} + \sqrt{1270} + \sqrt{172570})^2} = \frac{173842}{(1.470 + 35.637 + 415.42)^2} \approx \frac{173842}{204853} \approx 0.849\]

\textbf{Down-type triplet (d, s, b):}

\[Q_{dsb} = \frac{4.67 + 93.4 + 4180}{(\sqrt{4.67} + \sqrt{93.4} + \sqrt{4180})^2} = \frac{4278.1}{(2.161 + 9.664 + 64.653)^2} \approx \frac{4278.1}{5849} \approx 0.731\]

Both values are substantially above \(2/3 = 0.6667\). The up triplet
deviates by approximately \(+27\%\) and the down triplet by
approximately \(+10\%\), far outside any reasonable measurement
uncertainty for these central values (top quark mass uncertainty is
\(\pm 0.3\) GeV, contributing \(\ll 0.001\) to \(Q_{uct}\)).

\textbf{Interpretation.} The lepton triplet satisfies the Koide formula
exactly; the quark triplets do not. This could be explained in several
ways. (1) The Koide relation is a property of the lepton sector only,
not a universal tholonic prediction. This restricts the tholonic claim
but does not falsify it: the structural derivation would need to
identify why leptons satisfy \(Q = 2/3\) and quarks do not. (2) The
\(\overline{\text{MS}}\) scheme introduces renormalization-group running
that depends on the strong coupling constant; evaluating all quark
masses at a single unified scale might give different values. (3) The
Koide form for quarks may require a modified tholonic mapping that
accounts for color charge. These are open questions. The computation is
reported accurately here so the constraint can be assessed honestly
rather than obscured by appeals to quark mass uncertainty.

\textbf{8.2 Virial Ratio Drift in Strongly Correlated Systems}

For pure Coulomb systems, the virial theorem gives
\(\langle T \rangle / |\langle V \rangle| = 1/2\). The tholonic
attractor \(\phi\) provides a candidate alternative balance ratio of
\(1/\phi \approx 0.618\). In highly correlated electron systems (Wigner
crystals, heavy atoms with large correlation energy fractions, strongly
correlated transition metal compounds), the effective virial ratio might
drift toward 0.618 if the tholonic attractor is operative. Using
existing high-accuracy quantum chemistry databases (CCSD(T) or DMRG
energies for a large sample of atoms and molecules), plot the effective
virial ratio against the correlation energy fraction. A monotonic trend
from 0.5 toward 0.618 with increasing correlation would be a
quantitative tholonic signature distinguishable from standard QM
expectations.

\textbf{8.3 Nuclear Binding Energy and the 2-3-6 Recursion}

The nuclear magic numbers (2, 8, 20, 28, 50, 82, 126) mark positions of
exceptional nuclear stability. The shell model derives them from
spin-orbit coupling with fitted parameters. Compute the sequence
generated by the 2-3-6 tholonic recursion and its products at successive
levels, and evaluate how well that sequence predicts the observed magic
numbers compared to the shell model. If the 2-3-6 recursion generates
the magic number sequence with fewer free parameters than the shell
model, that is a structural prediction of tholonic stability with
measurable predictive economy.

\subsubsection{Category 2: Requiring New
Experiments}\label{category-2-requiring-new-experiments}

\textbf{8.4 Bell Violation Magnitude at Tetrahedral Measurement Angles}

The Tsirelson bound (\(2\sqrt{2}\)) is the universal maximum for CHSH
violation. Standard QM predicts the maximum at specific detector angle
pairs but does not predict structure in sub-maximum violations as a
function of the full measurement-axis geometry. If tholonic geometry
predicts specific local enhancements of Bell violation when detector
axes are oriented at tetrahedral angles (\(\arccos(1/3) \approx 70.53^\circ\)
relative to each other), this is distinguishable from standard QM. Bell
test platforms using entangled photons or trapped ions with full angular
control over detector orientations could scan violation magnitude as a
function of measurement-axis angle and look for structure at tetrahedral
positions.

\textbf{8.5 Quasicrystal Stability at \(\phi\)-Ratio Compositions}

The tholonic model predicts \(\phi\) as the convergence attractor of
optimal D-C balance. Quasicrystals with icosahedral symmetry have
inherent \(\phi\)-related geometry, but component atom ratios vary
continuously across alloy families. The prediction is that alloys whose
component atomic ratios are nearest to \(\phi \approx 1.618\) or
\(1/\phi \approx 0.618\) should show anomalously high thermal stability
(higher melting temperature, lower phason strain, longer structural
coherence length) compared to neighboring compositions. A systematic
synthesis-and-measurement campaign across a well-chosen alloy family
(e.g., Al-Mn-Si or Al-Cu-Fe icosahedral systems) varying composition in
fine steps around the \(\phi\) ratio would test this directly.

\subsubsection{Category 3: Requiring Theoretical Derivation
Work}\label{category-3-requiring-theoretical-derivation-work}

\textbf{8.6 Fine Structure Constant from Thologram Geometry}

The most consequential open derivation. The fine structure constant
(CODATA 2022: \(1/\alpha = 137.035999206\)) is not derived from any
deeper theory in the Standard Model; it is measured and inserted.
Several notable physicists have attempted expressions for \(1/\alpha\):
Arthur Eddington famously argued on numerological grounds that
\(1/\alpha\) is exactly 137 (an integer), and Richard Feynman described
\(\alpha \approx 1/137\) as ``one of the greatest damn mysteries of
physics.'' These prior attempts are cautionary context. Eddington's
derivation was rejected not because the numerical coincidence was
uninteresting but because his justification had no mechanism: he chose a
structure to fit the number rather than deriving the number from the
structure. The tholonic approach avoids this by constraining the search
to expressions whose structural elements (constants, exponents) are
independently fixed by the thologram geometry before any comparison with
the measured value is made. A systematic computational search over
tholonic-constant expressions has identified a leading candidate and a
secondary one.

\textbf{Leading candidate:
\(1/\alpha \approx 4 \cdot e^5 \cdot (\ln 2)^4 = 137.035865\), error
\(-0.98\) ppm.}

Since \((\sqrt{2})^4 = 2^2 = 4\) exactly, this is equivalently
\(e^5 \cdot (\sqrt{2})^4 \cdot (\ln 2)^4\), a product of three tholonic
constants with integer exponents. Its tholonic reading is:

\[\frac{1}{\alpha} \approx C_\text{outer} \cdot e^{D_\text{axis}} \cdot (\ln 2)^{C_\text{outer}}\]

where \(C_\text{outer} = 4\) is the C outer vertex value (\(2^2\)),
\(D_\text{axis} = 5\) is the Definition axis multiplier, and the same
\(C_\text{outer} = 4\) reappears as the exponent of \(\ln 2\). All three
structural numbers are fixed thologram values, not fitted parameters. No
alternating series with tholonic seeds, no two-constant product with
integer exponents up to \(\pm 8\), and no four-constant product within
the searched range came within \(200\) ppm. This candidate is the unique
expression in the three-tholonic-constant space within 1 ppm of
\(1/\alpha\).

The formula's structure encodes the thologram's D-C relationship
directly: the D axis multiplier (5) governs the exponential growth rate
via \(e\), and the C outer vertex value (4) appears twice, as both the
scale factor and the power of the logarithmic decay via \(\ln 2\). The C
role in the thologram is the accumulating, integrating component;
\(\ln 2\) is the tholonic constant associated with logarithmic
(integrating) convergence. The D role is the bounding, limiting
component; \(e\) raised to the D-axis power sets the overall magnitude.
Whether this structural correspondence can be derived from role axioms
rather than observed numerically is the open theoretical question.

\emph{Secondary candidate:
\(1/\alpha \approx 4\pi^3 + \pi^2 + \pi = 137.036304\), error \(+2.22\)
ppm.} This polynomial in \(\pi\) with tholonic-integer coefficients can
be written \(C \cdot \pi^{D+N} + N \cdot \pi^D + N \cdot \pi^N\) (outer
vertices \(N=1, D=2, C=4\)), but it is 2.3 times less accurate than the
leading candidate and the D vertex is structurally absent from the
prefactors, which requires justification. Note that
\((A + B)/2 = 137.036084\) at \(+0.62\) ppm is closer than A alone but
further than B alone; a weighted average has no tholonic justification
and is not a derivation.

Neither candidate is exact at the measurement precision of \(\alpha\)
(\(0.15\) ppb). The leading candidate is nonetheless the
best-constrained numerical target available: a tholonic derivation of
\(\alpha\) should produce an expression in the neighborhood of
\(C_\text{outer} \cdot e^{D_\text{axis}} \cdot (\ln 2)^{C_\text{outer}}\)
or explain why it coincides to 1 ppm.

\textbf{8.7 Proton/Electron Mass Ratio from N-D-C Weights}

The ratio \(m_p/m_e \approx 1836.15\) is not derived in the Standard
Model. The approximation \(6\pi^5 \approx 1836.12\) is known empirically
but unexplained. In the tholonic hydrogen atom, the proton is the D-role
component and the electron is the C-role component, with the N weight (6
in the 2-3-6 pattern) as the mediating value. The target derivation is:
does \(6\pi^5\) emerge naturally from the tholonic assignment (N = 6,
\(\pi\) from the tholonic constant set) under some specific structural
operation on the D and C weights? This is speculative, but the known
approximation \(6\pi^5 \approx 1836\) constrains the form the derivation
must take, making it a structured search rather than an open-ended one.

\textbf{8.8 CHSH Derivation (Complete at the Tholonic Structural Level)}

As detailed in Section 3.4, the tholonic derivation of the CHSH operator
and Tsirelson bound is now complete at the structural level. The value
\(2\sqrt{2}\) follows from two independent tholonic arguments: (a)
real-virtual tholon pair convergence (each component converges to
\(\sqrt{2}\), composite = \(2\sqrt{2}\)), and (b) pure angle geometry:
the optimal CHSH measurement angle is \(\pi/4\) (the tholonic D-C
balance angle and the \(\pi\)-branch seed from Paper 1),
\(\cos(\pi/4) = 1/\sqrt{2}\) is the tholonic convergence limit, and
\(\mathrm{CHSH}_{\max} = 4\cos(\pi/4) = 2\sqrt{2}\). The sign pattern is
derived from T1+T2+T3, where T3 (D\(\times\)C balance sign rule) is
itself derived from Paper 3's Lemma 4.2 and Definition 4.3 without
reference to measurement scenarios. A companion computational script
(\texttt{scripts/chsh\_tholon\_uniqueness.py}) exhaustively verifies all
claims.

The remaining item outside the tholonic framework is Bell's theorem
itself: WHY quantum systems achieve \(\cos\theta\) correlations while
classical systems are limited to 2. This is not a gap in the tholonic
argument; it is the boundary of what tholonic axioms address. The
tholonic derivation is a structural account of CHSH, not a re-derivation
of quantum mechanics from scratch.

\subsubsection{Category 4: Questions the Standard Model Cannot
Address}\label{category-4-questions-the-standard-model-cannot-address}

These are not experimental tests. They are foundational questions where
the Standard Model is structurally silent and the tholonic model
provides at least a principled structural answer.

\textbf{Why exactly three spatial dimensions?} The tetrahedron is the
minimum structure that encloses a volume. In fewer than three spatial
dimensions, no volume is enclosed and no thologram can form. In more
than three dimensions, the tetrahedron is no longer the minimal
enclosing structure and tholonic self-similarity is not uniquely
defined. If the tholonic model formally proves that stable recursive
N-D-C structure requires exactly three spatial dimensions, this answers
a question the Standard Model cannot frame.

\textbf{Why exactly three fermion generations?} Paper 3 proves three
roles are the minimum for non-trivial recursive stability. If this
requirement applies at the generation level, three generations is a
structural prediction. No fourth-generation fermion has been found,
which is consistent with this prediction but not specifically predicted
by standard QM.

\textbf{Why is the universe mathematically comprehensible?} The Standard
Model does not address this question. The tholonic model, with A\&I as
ontologically primary, makes comprehensibility a structural consequence:
a universe generated by Awareness and Intention is, by construction, one
whose patterns are legible to awareness. This is not a testable
prediction but it is an answer to a foundational question that physics
otherwise cannot frame.

\begin{center}\rule{0.5\linewidth}{0.5pt}\end{center}

\section{Implications}\label{implications}

Tesla's and Einstein's objections to the electron-as-discrete-particle
are documented, serious, and partially vindicated by quantum field
theory. The tholonic model addresses their concern by providing a
specific structural grammar for the field-first ontology both sought:
particles are stable tholonic excitations of the A\&I field, and their
measured properties are tholonic coordinates. The unified field that
Einstein did not find in geometry is proposed to be the A\&I field of
the tholonic model.

The atom's status as a measurable tholon currently rests on structural
correspondences that the tholonic framework generates from its own
geometry and that measurement independently confirms. The Koide
formula's \(Q = 2/3\) is currently the strongest case: the Standard
Model has no derivation of this value, while the tholonic D:C weight
ratio generates 2/3 directly from the 2-3-6 structural weights.

The case that the tholonic model is a \emph{better} model than the
Standard Model, not merely a consistent one, is argued in Section 6. The
Standard Model's 17 postulated fields and 19 measured parameters
represent an explanatory gap that the tholonic model proposes to fill by
reducing axioms and deriving parameters structurally. The transition
from structural account to predictive theory requires completing at
least one of the quantitative derivations in Section 8. The fine
structure constant (\(\alpha \approx 1/137\)) remains the priority
target: a tholonic derivation of \(\alpha\) from thologram geometry
would constitute the single strongest available confirmation of the
framework.

\begin{center}\rule{0.5\linewidth}{0.5pt}\end{center}

\section{References}\label{references-and-source-materials}

\begin{itemize}
\tightlist
\item
  \textbf{Tholonia: The Mechanics of Existential Awareness} by Duncan
  Stroud --- especially Chapters 11 (Tholon), 12 (Thologram), 13
  (Instances), 14 (Applications), and 15 (Mind).
\item
  \textbf{Paper 1 in this series:} \emph{Emergence of Classical
  Constants from a Minimal Recursive Triadic Framework} --- derivation
  of \(\pi/4\), \(\phi\), \(e\), \(\sqrt{2}\), and \(\ln 2\) from
  tholonic recursion.
\item
  \textbf{Paper 3 in this series:} \emph{Minimal Recursive Triadic
  Framework} --- formal proof that three functional roles are the
  minimum for non-trivial convergent recursion.
\item
  \textbf{Koide, Y.} (1982). ``A fermion-boson composite model of quarks
  and leptons.'' \emph{Physics Letters B} 120, 161--165. (Koide
  formula.)
\item
  \textbf{Einstein, A., Podolsky, B., Rosen, N.} (1935). ``Can
  Quantum-Mechanical Description of Physical Reality Be Considered
  Complete?'' \emph{Physical Review} 47, 777.
\item
  \textbf{Einstein, A.} (1949). ``Autobiographical Notes.'' In
  \emph{Albert Einstein: Philosopher-Scientist}, ed.~P.A. Schilpp. Open
  Court.
\item
  \textbf{Cirel'son, B.S.} (1980). ``Quantum generalizations of Bell's
  inequality.'' \emph{Letters in Mathematical Physics} 4, 93--100.
  (Tsirelson bound.)
\item
  \textbf{Aspect, A., Grangier, P., Roger, G.} (1982). ``Experimental
  Realization of Einstein-Podolsky-Rosen-Bohm Gedankenexperiment: A New
  Violation of Bell's Inequalities.'' \emph{Physical Review Letters} 49,
  91.
\item
  \textbf{Griffiths, D.J.} (2005). \emph{Introduction to Quantum
  Mechanics}, 2nd ed.~Pearson.
\item
  \textbf{Weinberg, S.} (1995). \emph{The Quantum Theory of Fields},
  Vol. 1. Cambridge University Press.
\item
  \textbf{Clausius, R.} (1870). ``On a Mechanical Theorem Applicable to
  Heat.'' \emph{Philosophical Magazine} 40, 122. (Virial theorem.)
\item
  \textbf{Particle Data Group.} (2024). \emph{Review of Particle
  Physics.} Progress of Theoretical and Experimental Physics 2024,
  083C01. (PDG quark and lepton mass values.)
\end{itemize}

\end{document}
